% \newcommand{\gtext}[1]{\text{\textcolor{ForestGreen}{#1}}}
% \newcommand{\leftleadsto}{\mathrel{\rotatebox[origin=c]{180}{$\leadsto$}}}
% \NewDocumentCommand{\typeable}{ o m m m }{
%   \IfValueTF{#1}
%     {  \mathrm{Typing}(\meta{#1(#2,#3,#4)}) }
%     {  \mathrm{Typing}(#2,#3,#4) }
% }
% \NewDocumentCommand{\typeable}{ o m m m }{
%   \IfValueTF{#1}
%     {  \meta{#1(#2)} \vdash \meta{#1(#3)}: \meta{#1(#4)} }
%     {  #2 \vdash #3: #4 }
% }
\newcommand{\stlctype}[3]{{#1} \vdash {#2} : {#3}}
\newcommand{\sketch}[1]{#1^{\circ}} 
\newcommand{\meta}[1]{
  \ensuremath{%
    #1
  }
}
\newcommand{\smeta}[1]{\ensuremath{%
    \setlength{\fboxsep}{0.1em}% 调整这个值来控制 padding
    \colorbox{gray!20}{
      \!$\displaystyle #1$\!
}}}
\newcommand{\acquire}[1]{\mathrm{Acquire}(#1)}
\NewDocumentCommand{\unify}{ m }{\!\mathrm{Unify}(\smash{#1})}
\newcommand{\ensure}[1]{#1} % 用于确保类型正确
\definecolor{skunknown}{RGB}{180,0,180}
\newcommand{\sk}[1]{\textcolor{skunknown}{#1}} 
\newcommand{\skhzc}[1]{#1}
\newcommand{\mvhzc}{\text{FV}}
\newcommand{\welltyped}[1]{\mathit{well\text{-}typed}\left(#1\right)}
% ===== 符号格式化命令 =====
\newcommand{\var}[1]{\mathit{#1}}           % 变量名称, 如 x, y
\newcommand{\type}[1]{\mathsf{#1}}          % 类型名称, 如 Int, Bool
\newcommand{\lang}[1]{\mathbf{#1}}          % 语言名称, 如 ML, Java
\newcommand{\tens}[1]{\mathbf{#1}}          % 张量名称, 如 T, U
\newcommand{\thm}[1]{\mathbf{#1}}           % 定理名称, 如 Thm1, Lemma2
\newcommand{\kw}[1]{\mathtt{#1}}            % 关键字, 如 if, while
\newcommand{\op}[1]{\mathrm{#1}}            % 操作符, 如 +, &&, ||
\newcommand{\lit}[1]{\mathrm{#1}}           % 字面量, 如 true, false
\newcommand{\func}[1]{\mathrm{#1}}          % 函数名称, 如 main, add
\newcommand{\set}[1]{\mathbb{#1}}           % 数学集合, 如 N, Z, R
\newcommand{\cat}[1]{\mathcal{#1}}          % 范畴, 如 C, D
% 自定义箭头命令
\newcommand{\updir}{\textcolor{black}{↑}}  % 可改为绿色 \textcolor{green}{↑}
\newcommand{\downdir}{\textcolor{black}{↓}}
%=============================================================================

\newcommand{\nbc}[3]{
	{\colorbox{#3}{\bfseries\sffamily\scriptsize\textcolor{white}{#1}}}
	{\textcolor{#3}{\sf\small \textit{#2}}}
}
\definecolor{rycolor}{RGB}{150, 0, 0}
\newcommand\ruyi[1]{\nbc{RJ}{#1}{rycolor}}
\newcommand\yx[1]{\nbc{YX}{#1}{rycolor}}
\newcommand{\toolname}[1]{\textsc{#1}} 
\newcommand{\tcode}[1]{{\small \texttt{\spaceskip=0.2em #1}}}
\newcommand{\scode}[1]{{\fontsize{8.2pt}{\baselineskip}\tt \texttt{\spaceskip=0.15em #1}}}
\newcommand{\redbox}[1]{{\setlength{\fboxsep}{2pt}\colorbox{red!20}{#1}}}
\newcommand{\greenbox}[1]{{\setlength{\fboxsep}{2pt}\colorbox{green!20}{#1}}}
\newcommand{\mainname}{\toolname{TyFlow}\xspace}
\newcommand{\stlc}{\lambda^{\rightarrow}}

\newcommand{\emphasis}[1]{\textcolor{rycolor}{#1}}

\def\modify#1#2#3{{\small\underline{\sf{#1}}:} {\color{red}{\small #2}}
{{\color{red}\mbox{$\Rightarrow$}}} {\color{blue}{#3}}}
\newcommand{\yxmodify}[2]{\modify{Yingfei}{#1}{#2}}
\newcommand{\ruyimodify}[2]{\modify{Ruyi}{#1}{#2}}
\newcommand{\ym}[2]{\yxmodify{#1}{#2}}
\newcommand{\ymok}[2]{{#2}}
\newcommand{\ymno}[2]{{#1}}